% Figure IV.8 --- the custody bound is the absence of an extra transition;
% the theorem is tested against its own bypass counterfactual.
%
% Reader question: what exactly does "non-bypassable custody" buy, and what
%   breaks the bound?
% Claim: the bound holds only because the custody ledger exposes exactly one
%   terminal transition; a single extra operator-controlled edge changes the
%   worst case from the agreed fee to the full balance.
% Evidence: Design Invariant (thm:fh-escrow-bound)'s four hypotheses; the two
%   terminal outcomes CLEAR/REFUSE; the extra redirect/second-close edge.
% Must distinguish: the one-edge difference between the two panels from
%   everything else, which is held identical.
% Grammar chosen: two panels on one grid (tpl-before-after.tex idiom), applied
%   to a hypothesis-present/hypothesis-absent counterfactual rather than a
%   time change -- the point is topological (one extra edge), so identical
%   geometry carries it. Grammar rejected: a single diagram with a dashed
%   "optional" edge would blur which panel obeys the theorem's hypothesis; a
%   table would lose the shape of the extra edge.
% Five-second test: which single edge is present in B but not A, and what
%   does it change the worst case to?
\pdfigurehue{pdgold}%
\begin{figure}[H]
\centering
\begin{tikzpicture}[pd figure,x=1cm,y=1cm]
  \def\ax{0.00}   \def\bx{4.95}
  \node[pd title,anchor=west,text width=4.6cm,align=left] at (\ax,2.75) {A. non-bypassable custody};
  \node[pd label,text=pdinkmuted,anchor=north west,text width=5.3cm,align=left] at (\ax,2.20) {one terminal transition};
  \node[pd title,anchor=west,text width=4.6cm,align=left] at (\bx,2.75) {B. operator bypass exists};
  \node[pd label,text=pdinkmuted,anchor=north west,text width=3.9cm,align=left] at (\bx,2.20) {same ceremony, weaker boundary};

  % ---- panel A: funded -> criteria -> Clear / Refuse ---------------------
  \node[pd datum] (fundeda) at (\ax+0.30,0) {};
  \node[pd label,anchor=north] at (\ax+0.30,-0.28) {funded};
  \node[draw=pdinkmuted!70,line width=.9pt,fill=pdpage,diamond,
    minimum width=1.55cm,minimum height=0.95cm,inner sep=0pt,pd label]
    (gatea) at (\ax+1.55,0) {criteria};
  \node[pd focus datum] (cleara) at (\ax+2.85,0.95) {};
  \node[pd datum] (refusea) at (\ax+2.85,-0.95) {};
  \draw[pd arrow] (fundeda) -- (gatea);
  \draw[pd focus arrow] (gatea) -- (cleara);
  \draw[pd arrow] (gatea) -- (refusea);
  \node[pd focus label,anchor=west] at (\ax+3.10,0.95) {\textbf{Clear}};
  \node[pd label,anchor=west] at (\ax+3.10,-0.95) {\textbf{Refuse}};
  \draw[pd focus rule] (\ax+0.05,-1.65) -- (\ax+4.55,-1.65);
  \node[pd focus label,anchor=north west,text width=4.5cm,align=left] at (\ax+0.05,-1.85)
    {worst-case loss $\leq\varphi$};

  % ---- panel B: identical graph, plus one operator-controlled edge -------
  \node[pd datum] (fundedb) at (\bx+0.30,0) {};
  \node[pd label,anchor=north] at (\bx+0.30,-0.28) {funded};
  \node[draw=pdinkmuted!70,line width=.9pt,fill=pdpage,diamond,
    minimum width=1.55cm,minimum height=0.95cm,inner sep=0pt,pd label]
    (gateb) at (\bx+1.55,0) {criteria};
  \node[pd focus datum] (clearb) at (\bx+2.85,0.95) {};
  \node[pd datum] (refuseb) at (\bx+2.85,-0.95) {};
  \node[pd breach datum] (lossb) at (\bx+2.85,-2.55) {};
  \draw[pd arrow] (fundedb) -- (gateb);
  \draw[pd focus arrow] (gateb) -- (clearb);
  \draw[pd arrow] (gateb) -- (refuseb);
  \draw[pd breach arrow,line width=1.25pt] (fundedb) to[out=-75,in=150] (lossb);
  \node[pd focus label,anchor=west] at (\bx+3.10,0.95) {\textbf{Clear}};
  \node[pd label,anchor=west] at (\bx+3.10,-0.95) {\textbf{Refuse}};
  \node[pd breach label,anchor=north,text width=4.2cm,align=center] at (\bx+2.85,-2.95)
    {redirect / second close: full balance exposed};
  \draw[pd breach rule] (\bx+0.05,-4.35) -- (\bx+4.55,-4.35);
  \node[pd breach label,anchor=north west,text width=3.9cm,align=left] at (\bx+0.05,-4.55)
    {worst case $=$ full custody balance};

  \node[pd note,anchor=north west,text width=10.6cm] at (\ax,-6.10)
    {the custody assumption: whitelist, fee cap, one terminal transition are
    the theorem's hypothesis, not decoration};
\end{tikzpicture}
\caption{\textbf{The custody bound is the absence of an extra transition.}
Panel A permits exactly one terminal move from funded custody to \textbf{Clear}
or \textbf{Refuse}: pay the whitelisted beneficiary or return the bond. With a
recipient whitelist and fee cap, worst-case loss is bounded by the agreed fee
$\varphi$. Panel B holds the same ceremony but drops the non-bypassability
hypothesis: one operator-controlled edge reaches redirect or a second close,
and the worst case becomes the full custody balance. Threshold signing
(\S\ref{sec:fh-escrow}) is only one candidate way to make panel A's absence of
that edge structural; the bound holds only when it actually is. [internal]}
\label{fig:fh-settlement}
\end{figure}
